%----------------------------------------------------------------------------------------
% RIJEN: BEGRIPPEN. Alle opgaven zijn eigen werk.
%
% Opbouw: Start, leerstof (kern, uitgewerkt voorbeeld, korte oefening), en daarna op een
% nieuwe pagina de Opdrachten (door elkaar, stijgende moeilijkheid), Koppel en Test jezelf.
% Labels van dit hoofdstuk beginnen met h1-.
%----------------------------------------------------------------------------------------

\mychapter{1}{Rijen: begrippen}

\lkr{Leerplan B+S: LPD 28 (recursief en expliciet voorschrift, rij als discrete functie). Dit hoofdstuk legt de begrippen vast voor hoofdstuk 2; rekenkundige en meetkundige rijen krijgen hier nog geen naam. Drie rode draden: het verschil tussen volgnummer en term, het verschil tussen expliciet en recursief, en de grafiek met losse punten.}

Het aantal inwoners van een stad elk jaar op 1 januari, het saldo van een spaarrekening elke maand, de blokjes in een trap die telkens een treetje hoger wordt: telkens krijg je getallen die na elkaar komen, in een vaste volgorde. Zo'n opeenvolging heet een rij. In dit hoofdstuk leer je hoe je een rij beschrijft: met een formule die elke term rechtstreeks geeft, of met een regel die zegt hoe je van de ene term naar de volgende gaat. Je kent al functies, en een rij blijkt er een bijzonder soort van te zijn. Die begrippen heb je nodig voor de rest van de cursus.

\begin{oefeningen}[Start]
\oef{h1-start-1}

\oef{h1-verken}
\end{oefeningen}

\lkr{Laat oefening \ref{oef:h1-verken} maken vóór de uitleg. In b redeneren leerlingen recursief (de vorige trap plus een kolom), in c expliciet (rechtstreeks uit het nummer van de trap). Kom hierop terug bij \textit{Expliciet of recursief?}.}

%========================================================================================
\section{Termen en volgnummers}

In de verkenning telde je blokjes: trap 1, trap 2, trap 3, enzovoort. Elk getal had een plaats in de rij. Om over zo'n rij te kunnen praten, geven we elk getal een naam met een index die zegt op welke plaats het staat.

\begin{theorie}
Een \emph{rij} is een opeenvolging van getallen in een vaste volgorde. De getallen zijn de \emph{termen} van de rij. De index is het \emph{volgnummer}: $u_1$ is de eerste term, $u_2$ de tweede, en $\formule{u_n}$ de $n$-de term of \emph{algemene term}, met $n \in \mathbb{N}_0$.

\vspace{1mm}
{\renewcommand{\arraystretch}{1.25}%
\begin{tabular}{@{}l*{7}{>{\centering\arraybackslash}p{11mm}}@{}}
volgnummer & 1 & 2 & 3 & $\cdots$ & $n-1$ & $n$ & $n+1$ \\
term & $u_1$ & $u_2$ & $u_3$ & $\cdots$ & $u_{n-1}$ & $u_n$ & $u_{n+1}$
\end{tabular}}
\end{theorie}

Het volgnummer zegt \textit{waar} een term staat, de term zelf zegt \textit{welk getal} daar staat. Houd die twee goed uit elkaar: $u_{n-1}$ is de term vóór $u_n$, en dat is iets anders dan $u_n - 1$. Drie puntjes na een rij betekenen dat ze verdergaat.

\voorbeeld{}{Het aantal dagen per maand in een gewoon jaar geeft een rij:\\[.5ex]
$31,\ 28,\ 31,\ 30,\ 31,\ 30,\ 31,\ 31,\ 30,\ 31,\ 30,\ 31$\\[.5ex]
Geef $u_2$ en $u_7$. Voor welke volgnummers is $u_n = 30$?\label{vb:h1-maanden}}{%
\uitleg{Bij $u_2$ en $u_7$ is het volgnummer gegeven en zoek je het getal op die plaats. Bij $u_n = 30$ is het omgekeerd: het getal is gegeven, en je zoekt de plaatsen.}
\begin{stappen}
u_2 = 28 & Op de tweede plaats staat februari, met 28 dagen. \\
u_7 = 31 & Op de zevende plaats staat juli, met 31 dagen. \\
n = 4,\ 6,\ 9 \text{ of } 11 & Het getal 30 staat op vier plaatsen: april, juni, september en november. Eenzelfde getal mag dus meer dan eens voorkomen, maar bij een volgnummer hoort maar één term.
\end{stappen}}

\lkr{Deze rij heeft twaalf termen en geen handig voorschrift. Zo zien leerlingen meteen dat een rij meer is dan een formule.}

\nujij{In de rij $4,\ 7,\ 7,\ 1,\ 0,\ 7$: wat is $u_3$? Voor welke volgnummers is $u_n = 7$?}{$u_3 = 7$; $n = 2$, 3 of 6.}

\begin{oefeningen}
\oef{h1-k1}
\end{oefeningen}

%========================================================================================
\section{Expliciet voorschrift}

Een rij met twaalf termen kun je gewoon opschrijven. Maar een rij die blijft doorgaan, kun je nooit helemaal uitschrijven. Dan heb je een voorschrift nodig: een regel waarmee je elke term kunt vinden. De eenvoudigste soort regel werkt net zoals een functievoorschrift: je vult het volgnummer in en krijgt de term.

\subsection{Termen berekenen}

\begin{theorie}
Een \emph{expliciet voorschrift} geeft $u_n$ rechtstreeks als uitdrukking in het volgnummer $n$, zoals $\formule{u_n = n^2 - 5n + 7}$.
\end{theorie}

Je kunt zo elke term berekenen zonder de vorige termen te kennen, net zoals je bij een functie $f(4)$ berekent zonder eerst $f(3)$ te zoeken.

\voorbeeld{}{Gegeven: $u_n = n^2 - 5n + 7$. Bereken $u_1$, $u_2$, $u_3$ en $u_{10}$. Schrijf ook $u_{n+1}$ zo eenvoudig mogelijk.\label{vb:h1-termen}}{%
\uitleg{Een term berekenen is invullen: vervang $n$ door het volgnummer. Bij $u_{n+1}$ vul je geen getal in, maar de uitdrukking $n + 1$.}
\begin{stappen}
u_1 = 1 - 5 + 7 = 3 & We vervangen $n$ door 1. \\
u_2 = 4 - 10 + 7 = 1, \quad u_3 = 9 - 15 + 7 = 1 & Twee termen mogen gelijk zijn, net zoals een functie twee keer dezelfde waarde kan aannemen. \\
u_{10} = 100 - 50 + 7 = 57 & De termen $u_4$ tot $u_9$ hebben we daarvoor niet nodig. \\
u_{n+1} = (n+1)^2 - 5(n+1) + 7 & Overal waar $n$ staat, komt nu $n+1$, met haakjes. \\
\phantom{u_{n+1}} = n^2 - 3n + 3 & Na het uitwerken krijgen we de volgende term. Die is dus niet hetzelfde als $u_n + 1$.
\end{stappen}}

\lkr{Laat leerlingen $u_{n+1}$ en $u_n + 1$ naast elkaar uitrekenen voor $n = 2$: 1 en 2. De indexnotatie is voor veel leerlingen het eerste struikelblok.}

\nujij{Gegeven: $u_n = 3n - n^2$. Bereken $u_1$, $u_4$ en $u_{n-1}$.}{$u_1 = 2$; $u_4 = -4$; $u_{n-1} = 3(n-1) - (n-1)^2 = -n^2 + 5n - 4$.}

\begin{oefeningen}
\oef{h1-k2}
\end{oefeningen}

\subsection{Voorschrift opstellen}

Vaak ken je alleen de eerste termen, en zoek je zelf het voorschrift. Daarvoor bestaat geen vaste formule: je zoekt een patroon. Het helpt om de termen te vergelijken met rijen die je kent, zoals $n$, $n^2$ of $2^n$.

\voorbeeld{}{Bepaal een expliciet voorschrift van $\tfrac13,\ {-\tfrac25},\ \tfrac37,\ {-\tfrac49},\ \ldots$\label{vb:h1-opstellen}}{%
\uitleg{Er zitten drie dingen in deze termen: een teken dat wisselt, een teller en een noemer. We zoeken ze elk apart, en zetten ze daarna samen.}
\begin{stappen}
\text{teken: } (-1)^{n+1} & De termen zijn positief bij oneven $n$ en negatief bij even $n$. $(-1)^{n+1}$ geeft precies $1,\ {-1},\ 1,\ {-1},\ \ldots$ \\
\text{teller: } 1,\ 2,\ 3,\ 4 \ \Rightarrow \ n & De teller is telkens het volgnummer. \\
\text{noemer: } 3,\ 5,\ 7,\ 9 \ \Rightarrow \ 2n+1 & De noemer wordt telkens 2 groter: dat is het dubbele van $n$, plus 1. \\[1ex]
u_n = (-1)^{n+1} \cdot \dfrac{n}{2n+1} & We controleren met een term: $u_4 = (-1)^5 \cdot \dfrac{4}{9} = -\dfrac{4}{9}$, en dat klopt.
\end{stappen}}

\begin{werkwijze}
Zoek het verband tussen volgnummer en term.
\begin{itemize}
  \item Vergelijk de termen met $n$, met $n^2$ en met machten zoals $2^n$. Tel er zo nodig iets bij of trek er iets van af.
  \item Bij een breuk bekijk je teller en noemer apart.
  \item Wisselend teken: $(-1)^n$ geeft $-1,\ 1,\ {-1},\ 1,\ \ldots$ en $(-1)^{n+1}$ geeft $1,\ {-1},\ 1,\ {-1},\ \ldots$
  \item Controleer je voorschrift met een term die je nog niet gebruikt hebt.
\end{itemize}
\end{werkwijze}

\thuis{\blokkop{Half uitgewerkt}\quad Rij $3,\ 6,\ 11,\ 18,\ \ldots$ Vergelijk met de kwadraten $1,\ 4,\ 9,\ 16$. Schrijf over in je schrift en vul aan: \ elke term is $\vak$ meer dan $n^2$, dus $u_n = \vak$ en $u_5 = \vak$. \quad {\footnotesize(Antwoord: 2 meer, $u_n = n^2 + 2$ en $u_5 = 27$.)}}

\nujij{Bepaal een expliciet voorschrift van $1,\ {-\tfrac12},\ \tfrac13,\ {-\tfrac14},\ \ldots$}{$u_n = (-1)^{n+1} \cdot \tfrac1n$.}

\begin{oefeningen}
\oef{h1-k3}
\end{oefeningen}

\subsection{Term of geen term}

Met een expliciet voorschrift kun je ook omgekeerd werken. Je kent een getal, en je vraagt je af of het in de rij voorkomt, en op welke plaats. Dan is de term gegeven, en zoek je het volgnummer: je lost een vergelijking op.

\voorbeeld{}{Gegeven: $u_n = n^2 + 2n$. Is 48 een term van de rij? En 50?\label{vb:h1-isterm}}{%
\naast{\uitleg{We stellen $u_n$ gelijk aan het getal en lossen op naar $n$. Een volgnummer moet een natuurlijk getal zijn, verschillend van 0: alleen zo'n oplossing telt.}
\begin{stappenvrij}{4.6cm}{4.6cm}
n^2 + 2n = 48 & We stellen de term gelijk aan 48. \\
(n-6)(n+8) = 0 & Na het ontbinden vinden we $n = 6$ of $n = -8$. Alleen $n = 6$ is een volgnummer, dus 48 is de 6de term. \\
n^2 + 2n = 50 & Nu met 50. \\
n = -1 \pm \sqrt{51} & De oplossingen zijn $n \approx 6{,}14$ en $n \approx -8{,}14$. Geen van beide is een volgnummer, dus 50 is geen term.
\end{stappenvrij}}{\begin{tikzpicture}[x=.5cm,y=.055cm,font=\scriptsize]
  \draw[black!14,very thin,xstep=1,ystep=10] (0,0) grid (7.4,70);
  \draw[->,grijs] (0,0) -- (7.8,0) node[right,text=black] {$n$};
  \draw[->,grijs] (0,0) -- (0,74) node[above,text=black] {$u_n$};
  \foreach \t in {1,...,7}{\node[below,inner sep=1.5pt] at (\t,0) {$\t$};}
  \foreach \t in {20,40,60}{\node[left,inner sep=1.5pt] at (0,\t) {$\t$};}
  \draw[grijs,dashed] (0,48) -- (7.4,48) node[right,text=black,yshift=-3pt] {48};
  \draw[grijs,dashed] (0,50) -- (7.4,50) node[right,text=black,yshift=3pt] {50};
  \foreach \x/\y in {1/3,2/8,3/15,4/24,5/35,6/48,7/63} \fill[accent] (\x,\y) circle (1.9pt);
\end{tikzpicture}}}

De figuur toont waarom: 48 is precies $u_6$, maar 50 ligt tussen $u_6 = 48$ en $u_7 = 63$. Er is geen plaats in de rij waar 50 staat.

\begin{werkwijze}
Is $x$ een term van de rij? Los $u_n = x$ op. Is er een oplossing $n \in \mathbb{N}_0$, dan is $x$ de $n$-de term. Anders is $x$ geen term.
\end{werkwijze}

\nujij{Gegeven: $u_n = n^2 + n$. Is 90 een term van de rij? En 100?}{90 is de 9de term ($n = 9$ of $n = -10$). Bij 100 is $n \approx 9{,}51$, dus geen term.}

\begin{oefeningen}
\oef{h1-k4}
\end{oefeningen}

\subsection{Starten bij volgnummer 0}

Tot nu toe begon elke rij bij $u_1$. Bij iets dat in de tijd verandert, is het vaak handiger om te starten bij $u_0$: de beginwaarde, voor er een stap gebeurd is. Dan is $u_n$ de waarde na $n$ stappen, en is $u_1$ al de tweede term. Het voorschrift verandert daardoor ook.

\voorbeeld{}{Een lift stijgt: de hoogte in meter is elke seconde $12,\ 15,\ 18,\ \ldots$ Geef een expliciet voorschrift als de eerste term $u_1$ heet, en als hij $u_0$ heet. Welk volgnummer heeft de hoogte 30 in elk geval?\label{vb:h1-nul}}{%
\uitleg{De hoogte stijgt telkens met 3, dus het voorschrift bevat $3n$. Wat erbij komt, hangt af van het volgnummer waarmee we starten: we kiezen het zo dat de eerste term 12 is.}
\begin{stappen}
u_1 = 12: \quad u_n = 3n + 9 & Bij $n = 1$ moet er 12 uitkomen, en $3 \cdot 1 + 9 = 12$. \\
u_0 = 12: \quad u_n = 3n + 12 & Bij $n = 0$ moet er 12 uitkomen, en $3 \cdot 0 + 12 = 12$. \\
3n + 9 = 30 \ \Rightarrow \ n = 7 & Start de rij bij $u_1$, dan is 30 de term $u_7$. \\
3n + 12 = 30 \ \Rightarrow \ n = 6 & Start de rij bij $u_0$, dan is 30 de term $u_6$: de hoogte na 6 seconden.
\end{stappen}}

Hetzelfde getal krijgt dus een ander volgnummer. Lees daarom altijd goed bij welke term een rij begint.

\lkr{In hoofdstuk 2 komt dit terug bij lineaire en exponentiële groei, waar $u_0$ de beginwaarde is.}

\begin{oefeningen}
\oef{h1-k5}
\end{oefeningen}

%========================================================================================
\section{Recursief voorschrift}

In de verkenning vond je trap 10 uit trap 9: je telde er gewoon een kolom van 10 blokjes bij. Je had daar geen formule voor nodig, alleen de vorige trap. Zo'n regel, die zegt hoe je van een term naar de volgende gaat, is een tweede manier om een rij vast te leggen.

\subsection{Termen berekenen}

\begin{theorie}
Een \emph{recursief voorschrift} geeft $u_n$ uit een of meer vorige termen, zoals $\formule{u_n = 2u_{n-1} - 1}$ voor $n \geqslant 2$. Je hebt ook de \emph{begintermen} nodig, zoals $\formule{u_1 = 3}$: zonder die ligt de rij niet vast.
\end{theorie}

Lees $u_n = 2u_{n-1} - 1$ als: een term is het dubbele van de vorige term, min 1. Hetzelfde voorschrift kun je ook schrijven met de index een plaats verschoven: $u_{n+1} = 2u_n - 1$ voor $n \geqslant 1$.

\thuis{\textit{Recursief} komt van het Latijnse \textit{recurrere}, teruglopen: je grijpt telkens terug naar de vorige term.}

\voorbeeld{}{Gegeven: $u_1 = 3$ en $u_n = 2u_{n-1} - 1$ voor $n \geqslant 2$. Bereken $u_2$, $u_3$ en $u_4$.\label{vb:h1-rec}}{%
\naast{\uitleg{We kennen alleen $u_1$. Elke volgende term halen we uit de term ervoor, dus we rekenen ze een voor een uit.}
\begin{stappenvrij}{4.6cm}{4.8cm}
u_2 = 2u_1 - 1 = 5 & Bij $n = 2$ is de vorige term $u_1 = 3$, en $2 \cdot 3 - 1 = 5$. \\
u_3 = 2u_2 - 1 = 9 & Nu gebruiken we de uitkomst van de vorige stap: $2 \cdot 5 - 1 = 9$. \\
u_4 = 2u_3 - 1 = 17 & En opnieuw: $2 \cdot 9 - 1 = 17$.
\end{stappenvrij}}{\begin{tikzpicture}[x=.6cm,y=.12cm,font=\scriptsize,>={Stealth[length=1.4mm]}]
  \draw[black!14,very thin,xstep=1,ystep=5] (0,0) grid (4.4,20);
  \draw[->,grijs] (0,0) -- (4.8,0) node[right,text=black] {$n$};
  \draw[->,grijs] (0,0) -- (0,22) node[above,text=black] {$u_n$};
  \foreach \t in {1,...,4}{\node[below,inner sep=1.5pt] at (\t,0) {$\t$};}
  \foreach \t in {5,10,15,20}{\node[left,inner sep=1.5pt] at (0,\t) {$\t$};}
  \foreach \a/\b/\c/\d in {1/3/2/5,2/5/3/9,3/9/4/17}{\draw[->,accent!60,shorten <=2pt,shorten >=2pt] (\a,\b) to[bend left=25] (\c,\d);}
  \foreach \x/\y in {1/3,2/5,3/9,4/17} \fill[accent] (\x,\y) circle (1.9pt);
\end{tikzpicture}}}

Elke pijl in de figuur is één stap van het voorschrift. Om $u_{100}$ te vinden, moet je zo 99 stappen zetten: recursief rekenen gaat stap per stap.

\nujij{Gegeven: $u_1 = 10$ en $u_n = u_{n-1} - 2n$ voor $n \geqslant 2$. Bereken $u_2$, $u_3$ en $u_4$.}{$u_2 = 6$; $u_3 = 0$; $u_4 = -8$.}

\begin{oefeningen}
\oef{h1-k6}
\end{oefeningen}

\subsection{Twee begintermen}

Een recursief voorschrift kan ook twee vorige termen gebruiken. Dan heb je twee begintermen nodig, anders kun je de derde term niet berekenen. De bekendste zo'n rij is de \textit{rij van Fibonacci}, waarin elke term de som is van de twee vorige: $u_1 = u_2 = 1$ en $u_n = u_{n-1} + u_{n-2}$ voor $n \geqslant 3$. Die geeft $1,\ 1,\ 2,\ 3,\ 5,\ 8,\ 13,\ \ldots$ De rij duikt op in de natuur, bijvoorbeeld in het aantal spiralen van een zonnebloem.

\voorbeeld{}{Gegeven: $u_1 = 2$, $u_2 = 3$ en $u_n = u_{n-1} \cdot u_{n-2} - 1$ voor $n \geqslant 3$. Bereken $u_3$, $u_4$ en $u_5$.\label{vb:h1-twee}}{%
\uitleg{Voor elke term hebben we nu de twee termen ervoor nodig. Daarom begint de rij met twee gegeven termen, en berekenen we vanaf $u_3$.}
\begin{stappen}
u_3 = u_2 \cdot u_1 - 1 = 3 \cdot 2 - 1 = 5 & Bij $n = 3$ zijn de twee vorige termen $u_2$ en $u_1$. \\
u_4 = u_3 \cdot u_2 - 1 = 5 \cdot 3 - 1 = 14 & Bij $n = 4$ schuift alles een plaats op: nu zijn het $u_3$ en $u_2$. \\
u_5 = u_4 \cdot u_3 - 1 = 14 \cdot 5 - 1 = 69 & En bij $n = 5$ zijn het $u_4$ en $u_3$.
\end{stappen}}

\begin{oefeningen}
\oef{h1-k7}
\end{oefeningen}

\subsection{Voorschrift opstellen}

Een recursief voorschrift opstellen is vaak makkelijker dan een expliciet: je hoeft alleen te kijken wat er gebeurt tussen twee termen die naast elkaar staan.

\voorbeeld{}{Bepaal een recursief voorschrift van $1,\ 2,\ 4,\ 7,\ 11,\ 16,\ \ldots$\label{vb:h1-recopst}}{%
\uitleg{We kijken wat er telkens bijkomt. Is dat niet altijd hetzelfde, dan zoeken we hoe het afhangt van het volgnummer.}
\begin{stappen}
\text{erbij: } 1,\ 2,\ 3,\ 4,\ 5 & Wat erbij komt, wordt telkens 1 groter, dus het hangt af van $n$. \\
u_2 - u_1 = 1,\quad u_3 - u_2 = 2 & Bij de stap naar $u_2$ komt er 1 bij, bij de stap naar $u_3$ komt er 2 bij. Van $u_{n-1}$ naar $u_n$ komt er dus $n - 1$ bij, en niet $n$. \\
u_n = u_{n-1} + (n-1),\ u_1 = 1 & Dit geldt voor $n \geqslant 2$. We controleren met een term: $u_5 = 7 + 4 = 11$.
\end{stappen}}

\begin{werkwijze}
Kijk hoe je van een term naar de volgende gaat: wat komt erbij, waarmee vermenigvuldig je? Hangt de stap af van het volgnummer, schrijf ze dan met $n$ en controleer ze bij de eerste stap. Vergeet de begintermen niet.
\end{werkwijze}

\lkr{De fout $u_n = u_{n-1} + n$ is hier heel gewoon. Laat leerlingen die fout zelf ontdekken door de eerste stap te controleren: dan is $u_2 = 3$.}

\begin{oefeningen}
\oef{h1-k8}
\end{oefeningen}

\subsection{Expliciet of recursief?}

Je kent nu twee manieren om een rij vast te leggen. Ze hebben elk hun eigen sterkte.

\begin{tabular}{@{}>{\RaggedRight\arraybackslash}p{7cm}>{\RaggedRight\arraybackslash}p{7.6cm}@{}}
\textbf{Expliciet:} rechtstreeks uit het volgnummer. Handig om meteen een verre term te berekenen. & \textbf{Recursief:} stap per stap uit de vorige termen. Handig om te beschrijven hoe een rij ontstaat. \\[2mm]
\begin{tikzpicture}[>={Stealth[length=1.8mm]},font=\small]
  \node (a) at (0,0) {$100$};
  \node (b) at (2.4,0) {$u_{100}$};
  \draw[->,accent,thick] (a) -- (b);
\end{tikzpicture} &
\begin{tikzpicture}[>={Stealth[length=1.8mm]},font=\small]
  \node (a) at (0,0) {$u_1$};
  \node (b) at (1.1,0) {$u_2$};
  \node (c) at (2.2,0) {$u_3$};
  \node (d) at (3.3,0) {$\cdots$};
  \node (e) at (4.5,0) {$u_{100}$};
  \draw[->,accent,thick] (a) -- (b);
  \draw[->,accent,thick] (b) -- (c);
  \draw[->,accent,thick] (c) -- (d);
  \draw[->,accent,thick] (d) -- (e);
\end{tikzpicture}
\end{tabular}

\vspace{1mm}
Niet elke rij heeft een eenvoudig voorschrift, van welke soort ook. Voor de priemgetallen 2, 3, 5, 7, 11, \ldots\ kent niemand een eenvoudige formule. Heb je wel een expliciet voorschrift, dan kun je er een recursief uit halen: bereken wat er bijkomt tussen $u_{n-1}$ en $u_n$.

\voorbeeld{}{Gegeven: $u_n = n^2$. Bepaal een recursief voorschrift.\label{vb:h1-omzet}}{%
\uitleg{We zoeken wat er bijkomt van de ene term naar de volgende. Daarvoor schrijven we de vorige term ook met $n$, en trekken we af.}
\begin{stappen}
u_{n-1} = (n-1)^2 = n^2 - 2n + 1 & De vorige term vinden we door $n$ te vervangen door $n - 1$. \\
u_n - u_{n-1} = 2n - 1 & Het verschil zegt wat er bijkomt: $n^2 - (n^2 - 2n + 1) = 2n - 1$. \\
u_n = u_{n-1} + 2n - 1,\ u_1 = 1 & Dit geldt voor $n \geqslant 2$. We controleren: $u_2 = 1 + 3 = 4$ en $u_3 = 4 + 5 = 9$.
\end{stappen}}

\lkr{Kom terug op oefening \ref{oef:h1-verken}: welke redenering was recursief en welke expliciet? Het aantal blokjes van trap $n$ is $u_n = u_{n-1} + n$ met $u_1 = 1$, of $u_n = \tfrac{n(n+1)}{2}$.}

\begin{oefeningen}
\oef{h1-k9}
\end{oefeningen}

\thuis{\blokkop{Met GeoGebra (thuis)}\quad Controleer je rekenwerk met twee commando's.
\begin{itemize}
  \item Expliciet: \texttt{Rij(3n+1,n,1,10)} geeft $u_1$ tot en met $u_{10}$ van $u_n = 3n + 1$.
  \item Recursief: \texttt{IteratieLijst(2x+4,-2,9)} geeft $u_1 = -2$ en daarna nog 9 termen van $u_n = 2u_{n-1} + 4$. De $x$ staat voor de vorige term.
\end{itemize}}

%========================================================================================
\section{De grafiek van een rij}

Bij een expliciet voorschrift zoals $u_n = n^2 - 6n + 10$ denk je meteen aan de functie $f(x) = x^2 - 6x + 10$. Dat is geen toeval: ook een rij koppelt aan elk getal precies één ander getal. Het verschil zit in de getallen die je mag invullen.

\begin{theorie}
Een rij is een \emph{functie met domein} $\formule{\mathbb{N}_0}$: bij elk volgnummer hoort precies één term. Haar \emph{grafiek} bestaat uit de losse punten $\formule{(n, u_n)}$.
\end{theorie}

\naast{In de grafiek staan de volgnummers horizontaal en de termen verticaal. Verbind de punten niet, want tussen twee volgnummers is er niets: $u_{2{,}5}$ bestaat niet.

\vspace{1.5mm}
De figuur toont de rij $u_n = n^2 - 6n + 10$. De stippellijn is de functie $f(x) = x^2 - 6x + 10$, met domein $\mathbb{R}$. De punten van de rij liggen op die parabool, maar de parabool zelf hoort niet bij de rij.}{\begin{tikzpicture}[x=.55cm,y=.24cm,font=\scriptsize]
  \draw[black!14,very thin,xstep=1,ystep=2] (0,0) grid (6.4,11);
  \draw[->,grijs] (0,0) -- (6.8,0) node[right,text=black] {$n$};
  \draw[->,grijs] (0,0) -- (0,12) node[above,text=black] {$u_n$};
  \foreach \t in {1,...,6}{\node[below,inner sep=1.5pt] at (\t,0) {$\t$};}
  \foreach \t in {2,4,...,10}{\node[left,inner sep=1.5pt] at (0,\t) {$\t$};}
  \draw[grijs,dashed,domain=.4:6.15,samples=40] plot (\x,{\x*\x-6*\x+10});
  \foreach \x/\y in {1/5,2/2,3/1,4/2,5/5,6/10} \fill[accent] (\x,\y) circle (1.9pt);
\end{tikzpicture}}

\voorbeeld{}{Teken de grafiek van $u_n = (-1)^n \cdot \dfrac{12}{n}$ voor $n = 1$ tot en met 6.\label{vb:h1-grafiek}}{%
\naast{\uitleg{We berekenen eerst de termen in een tabel, en zetten ze daarna als losse punten uit.}

\vspace{1mm}
\begin{tabular}{@{}l*{6}{>{\centering\arraybackslash}p{7.5mm}}@{}}
$n$ & 1 & 2 & 3 & 4 & 5 & 6 \\
$u_n$ & $-12$ & $6$ & $-4$ & $3$ & $-2{,}4$ & $2$
\end{tabular}

\vspace{2mm}
Door de factor $(-1)^n$ wisselt het teken: de punten liggen om en om onder en boven de horizontale as. Door de breuk $\tfrac{12}{n}$ komen ze steeds dichter bij die as. Een lijn tussen de punten zou iets tonen dat er niet is.}{\begin{tikzpicture}[x=.55cm,y=.12cm,font=\scriptsize]
  \draw[black!14,very thin,xstep=1,ystep=2] (0,-13) grid (6.4,8);
  \draw[->,grijs] (0,0) -- (6.8,0) node[right,text=black] {$n$};
  \draw[->,grijs] (0,-13) -- (0,9) node[above,text=black] {$u_n$};
  \foreach \t in {1,3,5}{\node[above,inner sep=1.5pt] at (\t,0) {$\t$};}
  \foreach \t in {2,4,6}{\node[below,inner sep=1.5pt] at (\t,0) {$\t$};}
  \foreach \t in {-12,-8,-4,4}{\node[left,inner sep=1.5pt] at (0,\t) {$\t$};}
  \foreach \x/\y in {1/-12,2/6,3/-4,4/3,5/-2.4,6/2} \fill[accent] (\x,\y) circle (1.9pt);
\end{tikzpicture}}}

\begin{oefeningen}
\oef{h1-k10}
\end{oefeningen}

%----------------------------------------------------------------------------------------
% OPDRACHTEN: door elkaar. Niveau 1, dan 2, dan 3.
%----------------------------------------------------------------------------------------
\opdrachtenkop

\lkr{De opdrachten staan bewust door elkaar, met stijgende moeilijkheid. In deze versie staat bij elke opdracht het onderwerp, zodat je kunt kiezen.}

\begin{oefeningen}
\oef{h1-termen}

\oef{h1-expl}

\oef{h1-recterm}

\oef{h1-opstellen}

\oef{h1-plant}

\oef{h1-recopst}

\oef{h1-isterm}

\oef{h1-o1}

\oef{h1-grafiek}

\oef{h1-o2}

\oef{h1-subst}

\oef{h1-o3}

\oef{h1-o4}

\oef{h1-nuleen}

\oef{h1-o5}

\oef{h1-o6}

\oef{h1-o7}

\oef{h1-grafvs}

\oef{h1-o8}

\oef{h1-omzet}

\oef{h1-o9}

\oef{h1-o10}

\oef{h1-recexpl}

\oef{h1-o11}

\oef{h1-o12}

\oef{h1-o13}

\oef{h1-o14}

\oef{h1-o15}

\oef{h1-o16}

\oef{h1-o17}

\oef{h1-o18}

\oef{h1-o19}

\oef{h1-o20}

\oef{h1-o21}

\oef{h1-o22}
\end{oefeningen}

%----------------------------------------------------------------------------------------
% KOPPEL: verbanden met andere leerstof
%----------------------------------------------------------------------------------------
\opdrachtsectie{Koppel}

\lkr{Rationale functies, goniometrie en matrices: kijk na of ze al gezien zijn in jouw jaarplanning.}

\begin{oefeningen}
\oef{h1-koppel-1}

\oef{h1-koppel-2}

\oef{h1-koppel-3}
\end{oefeningen}

%----------------------------------------------------------------------------------------
\opdrachtsectie{Test jezelf}

\begin{oefeningen}
\oef{h1-t1}

\oef{h1-t2}

\oef{h1-t3}

\oef{h1-t4}

\oef{h1-t5}

\oef{h1-t6}

\oef{h1-t7}

\oef{h1-t8}

\oef{h1-t9}
\end{oefeningen}

\herhaalschema{%
  \fv{oef:h1-t1}{vb:h1-maanden}{oef:h1-termen}%
  \fv{oef:h1-t2}{vb:h1-termen}{oef:h1-expl}%
  \fv{oef:h1-t3}{vb:h1-opstellen}{oef:h1-opstellen}%
  \fv{oef:h1-t4}{vb:h1-isterm}{oef:h1-isterm}%
  \fv{oef:h1-t5}{vb:h1-twee}{oef:h1-recterm}%
  \fv{oef:h1-t6}{vb:h1-recopst}{oef:h1-recopst}%
  \fv{oef:h1-t7}{vb:h1-omzet}{oef:h1-omzet}%
  \fv{oef:h1-t8}{vb:h1-grafiek}{oef:h1-grafiek}%
  \fv{oef:h1-t9}{vb:h1-nul}{oef:h1-nuleen}}

\ifzelfstudie
\section*{Maak je eigen samenvatting}
Sluit dit boek en schrijf in je schrift: wat is een term, een volgnummer en de algemene term, wat is het verschil tussen een expliciet en een recursief voorschrift, hoe stel je ze op, hoe test je of een getal een term is, wanneer start je bij $u_0$, hoe ziet de grafiek van een rij eruit, en welke fouten wil je vermijden? Vergelijk daarna met de kaders.
\fi
